\subsection{Nonlinear repairs and a uniform global partial model}

The following constructions isolate the remaining integral comparison
problem. Neither construction asserts a completed strict comparison
over $\Z$: the boundary-relative simplicial lifts and the eventual
infinite arithmetic descent remain to be supplied.

\begin{theorem}[C18: \status{Proved}; a nonlinear local disk closure]
\label{sa08:C18}
Fix a prime $p$, put $R=\Z_{(p)}$, and let $r,w\geq1$.
The strict PD differential algebra
\[
 D=\Lambda_R(x)\otimes_R\Gamma_R(y),\qquad
 |x|=2r-1,\quad |y|=2r,\quad dx=y,
 \qquad \operatorname{wt}(x)=\operatorname{wt}(y)=w
\]
embeds in a free graded PD algebra $C$ with a nonlinear differential
such that $C\to R$ is a quasi-isomorphism after every base change from
$R$. Every new generator has cochain degree at least $2rp-2$.
The construction preserves $q\leq w$ whenever that support condition
holds for the initial generators.

After adjoining the coefficient algebra $S=\Gamma_R(T)$, the new
generators can also be delayed by arbitrary sufficiently large
coefficient-weight thresholds: there are free-term subalgebras
$S\otimes D\subset E\subset S\otimes C$ that still resolve $S$
vertically in every fixed cochain range in all sufficiently large
weights, universally under base change.
\end{theorem}
\begin{proof}
Set $y_i=\gamma_{p^i}(y)$ and
\[
 c_i=\frac{(p^{i+1})!}{((p^i)!)^p}=p a_i.
\]
Legendre's formula gives $v_p(c_i)=1$, so $a_i\in R^\times$.
Divided-power multiplication gives $y_i^p=p a_i y_{i+1}$.
Use odd generators $x_i$, $i\geq0$, with $x_0=x$, and even
generators $v_i$, $i\geq0$, and define
\[
 C=\Gamma_R(y)\otimes\Lambda_R(x_i:i\geq0)
       \otimes\Gamma_R(v_i:i\geq0),
\]
where
\begin{align*}
 |x_i|&=2rp^i-1,&\operatorname{wt}(x_i)&=p^iw,\\
 |v_i|&=2rp^{i+1}-2,&\operatorname{wt}(v_i)&=p^{i+1}w,\\
 dx_i&=y_i,&dv_i&=p a_i x_{i+1}-x_i y_i^{p-1},\qquad dy=0.
\end{align*}
Extend $d$ as a graded PD derivation. The relation
$d^2v_i=p a_i y_{i+1}-y_i^p=0$ verifies that it is a
differential. The displayed degrees and weights verify the grading
claims and the lower bound on new generators.

Modulo $p$, the digit presentation
\[
 \Gamma_{\mathbb F_p}(y)
 =\bigotimes_{i\geq0}\mathbb F_p[y_i]/(y_i^p)
\]
makes $C\otimes\mathbb F_p$ the tensor product of the complexes
\[
 \mathbb F_p[y_i]/(y_i^p)\otimes\Lambda(x_i)\otimes\Gamma(v_i),
 \qquad dx_i=y_i,\quad dv_i=-x_i y_i^{p-1}.
\]
Each is its constant copy of $\mathbb F_p$ plus the following
two-term identity complexes, with the indicated sign:
\begin{align*}
 x_i y_i^a\gamma_m(v_i)&\longmapsto
       y_i^{a+1}\gamma_m(v_i)
       &&(0\leq a\leq p-2,\ m\geq0),\\
 \gamma_{m+1}(v_i)&\longmapsto
       -x_i y_i^{p-1}\gamma_m(v_i)
       &&(m\geq0).
\end{align*}
They exhaust the nonconstant monomials without overlap. The extra
terms vanish respectively by $x_i^2=0$ and $y_i^p=0$.
Finite tensor products and then their filtered union therefore have
cohomology $\mathbb F_p$ in degree zero only.

Every $C^q$ is finite free over $R$, because the generator degrees
are positive and tend to infinity. Thus $H^q(C)$ is a finite
$R$-module. The universal-coefficient injection
\[
 H^q(C)/pH^q(C)\lhook\joinrel\longrightarrow
 H^q(C\otimes_R\mathbb F_p)
\]
and Nakayama's lemma give $H^q(C)=0$ for $q>0$; also $C^0=R$.
Cycles and boundaries are free over the PID $R$, so the exact
cycle sequences split. The splitting can be made weight by weight.
Consequently $C$ is $R$ plus a contractible complex of free modules,
and the assertion survives arbitrary base change.

For the delayed version give $y,x_0$ delay zero, choose arbitrary
nonnegative integers $\nu(x_i)$ for $i\geq1$, and choose
\[
 \nu(v_i)\geq\max\{\nu(x_i),\nu(x_{i+1})\}.
\]
For a divided-power/exterior basis monomial $m$, let $\nu(m)$ be
the sum of its generator delays with multiplicity. Define $E$ to
be the $R$-span of the monomials $T^{[s]}m$ with $s\geq\nu(m)$.
Delay is additive under products and scales under divided powers.
Neither replacement in $dv_i$ increases it, and $dx_i$ removes the
new generator. Thus $E$ is a PD differential subalgebra containing
$S\otimes D$, and its homogeneous terms are free.

For each fixed cochain degree there are finitely many $m$. If
$W\geq\max_m\{\operatorname{wt}(m)+\nu(m)\}$, every term of
$(S\otimes C)^{W,q}$ belongs to $E$. Take this finite maximum
through degree $d+1$, including all preceding degrees, to compare
cohomology through degree $d$. The complexes then agree in the
required degrees in every larger weight, proving the eventual
vertical comparison and its base-change assertion.
\end{proof}

\paragraph{What C18 does not supply.}
To extend a prescribed map $D\to A$, one must choose compatible
primitives of $\gamma_{p^i}(dx)$ and then of the odd cycles defining
$dv_i$. Eventual exactness only gives individual primitives in a
fixed cochain degree at sufficiently large weights. It does not
give all those primitives at the prescribed weights $p^iw$ for one
$w$. The delayed version replaces this numerical issue by the need
for $S$-linear choices on coefficient tails, compatible with the
simplicial operators. Those choices have not been constructed here.
An absolute contraction of a parameter module does not settle the
relative problem: a horizontal cell has a prescribed boundary and
a sphere as its new quotient parameter.

\begin{theorem}[C21: \status{Proved}; a common global partial model]
\label{sa08:C21}
Let $A,B$ be generalized strict PD differential $S$-algebras over a
commutative ring $k$, satisfying the uniform eventual horizontal
acyclicity condition of scan~119. There is a simplicial nonnegative
bigraded strict PD differential $S$-algebra $P$, with morphisms
$P\to A$ and $P\to B$, such that:
\begin{enumerate}
\item every homogeneous term of $P$ is free over $k$;
\item $P$ satisfies uniform eventual horizontal acyclicity;
\item the coefficient map identifies $H^0(P_n^{w,*})$ with $S^w$
      in every simplex and weight;
\item after base change to $k_{(p)}=k\otimes\Z_{(p)}$, the
      coefficient map is a quasi-isomorphism through cochain degree
      $d$ whenever $d<2p-1$, uniformly in every simplex and weight;
      if $k_{(p)}$ is torsionfree over $\Z_{(p)}$, $d<2p$ suffices;
\item through degree $d$ the same comparison is an isomorphism
      after inverting $\lfloor(d+1)/2\rfloor!$.
\end{enumerate}
The construction preserves $q\leq w$ if that condition is imposed.
No vertical exactness hypothesis on $A,B$ is needed. In particular,
the theorem does not assert that $P$ is an admissible theory over $k$.
\end{theorem}
\begin{proof}
Write $K(n)=\operatorname{DK}(k[n])$, and let $D(n+1)$ be the
Dold--Kan image of the normalized identity disk with $k$ in
homological degrees $n+1,n$. There is an inclusion
$K(n)\subset D(n+1)$ that splits degreewise with a free complement.
Its quotient has normalized type $k[n+1]$; the prescribed boundary
is retained in the following attachment.

For a simplicial module $M$ and a positive weight $w$, let
$F_{q,w}(M)$ be the free strict PD differential $S$-algebra on a
vertical identity disk with copies of $M$ in cochain degrees
$q,q+1$, both of weight $w$. Its universal property is
\[
 \Hom_{\mathrm{PD\text{-}DGA}}(F_{q,w}(M),X)
 =\Hom_{s\operatorname{Mod}_k}(M,X^{w,q}).
\]
A normalized cycle $z\in N_n(X^{w,q})$ therefore gives a
relative cell attachment
\[
 X\longrightarrow
 X\amalg_{F_{q,w}(K(n))}F_{q,w}(D(n+1)).
\]
The new normalized element has boundary $z$. In each simplex this
is a free extension by independent vertical disks: the degreewise
free complement in the parameter module becomes free algebra
generators in the two consecutive cochain degrees. Thus no old
algebraic relation is changed.

Suppose $X$ maps to $A$ and $B$, and both target rows of bidegree
$(w,q)$ are acyclic. Each image of $z$ has a normalized primitive.
Choosing the two primitives extends the parameter maps from $K(n)$
to $D(n+1)$. The free universal property extends both algebra maps;
the images in degree $q+1$ are the differentials of the chosen
degree-$q$ images. Only target acyclicity in row $q$ is used.

Start with $S$ and its two coefficient maps. Choose
\[
 W(q)\geq\max\{1,q+1,\mu_A(q),\mu_B(q)\}.
\]
At stage $q$, attach simultaneously a cell for every normalized
cycle in every row $X^{w,q}$ with $w\geq W(q)$ and every homological
degree $n\geq0$. Repeat countably many times. Every cycle in the
colimit occurs at a finite round and is killed at the next round;
exact filtered colimits commute with normalization. These rows are
therefore exact, including in homological degree zero.

Perform these stages for $q=0,1,2,\ldots$. A stage with $q>0$
adjoins no generator of smaller cochain degree and does not alter
the lower rows. The final colimit consequently has the uniform
horizontal bound
\[
 \mu_P(d)=\max_{0\leq q\leq d}W(q).
\]
All maps to $A,B$ extend throughout. Every simplex algebra is free
over $S$ on the independent vertical disks added at that simplex,
so all terms are free over $k$. Positive generator weights preserve
the $k$-augmentation. The optional inequality $q\leq w$ follows
from $w\geq q+1$ for each new disk; that part of the bound is
otherwise unnecessary.

It remains to compute the vertical range. An even-to-odd disk
$\Gamma(x)\otimes\Lambda(y)$, $dx=y$, is its constants plus
the identity pairs
\[
 \gamma_m(x)\longmapsto\gamma_{m-1}(x)y,
 \qquad m\geq1.
\]
An odd-to-even disk, with odd $q=|x|\geq1$, is its constants
plus the two-term complexes
\[
 x\gamma_{m-1}(y)\xrightarrow{\ m\ }\gamma_m(y),
 \qquad
 \text{in degrees }m(q+1)-1,\ m(q+1).
\]
Over $\Z_{(p)}$, every pair with $m<p$ contracts by a unit.
The first possible nonunit pair begins in degree
$p(q+1)-1\geq2p-1$. Thus after arbitrary base change the
coefficient comparison is exact through every $d<2p-1$.
For torsionfree coefficients the first multiplication by $p$ is
injective, improving this to $d<2p$. The distinction is necessary:
for general coefficients its kernel can already contribute in
degree $2p-1$.

For completeness, expand a finite tensor product as a direct sum
of tensors of these pairs and constants. A summand containing a
unit pair contracts. If every selected pair is nonunit, the summand
starts above the stated range; in the torsionfree endpoint case,
the only possible first-degree contribution is one injective
$p$-pair. This proves the tensor assertion without assuming
flatness of the resulting cohomology. Infinite tensor products
are exact filtered colimits of finite ones. Tensoring with $S$
and taking a weight component preserve the conclusion. The bound
does not depend on the number of generators, hence is uniform in
the simplex and weight.

If a pair meets degrees at most $d$, then
$m(q+1)-1\leq d$ and therefore
$m\leq\lfloor(d+1)/2\rfloor$. Inverting the stated factorial
makes every such pair a unit pair, proving the final assertion.
Only the degree-zero even-to-odd disks affect cochain degree zero;
their universal contractions also prove $H^0(P_n^{w,*})=S^w$
in every weight over every $k$.
\end{proof}

\paragraph{The remaining integral construction.}
For each prime $p$, choose a $p$-power $N_p\geq W(0)$.
The marked lift of Statement~\ref{sa08:C09} gives
$G_{N_p}\to P_{(p)}\to A_{(p)},B_{(p)}$ using only the
acyclic degree-zero row of $P$. These local sources still have
the nonuniform initial weights ruled out by
Statement~\ref{sa08:C12}. A possible next construction would
replace $P_{(p)}$ by an admissible local extension retaining its
maps to both targets and changing only cochain degrees above a
bound tending to infinity with $p$. Statement~\ref{sa08:C18}
establishes such a degree bound for one independent disk.
It does not establish the required boundary-relative simplicial
extensions. Even after those extensions, infinite-prime arithmetic
descent would require its own proof. Neither step is assumed here.
